\section{Implications for the hybrid-solver program}
\label{sec:implications}

The conversation leads to four firm conclusions.

\begin{enumerate}[leftmargin=2em]
\item \emph{Acceleration has not vanished spectrally.}
The spider's restricted condition number is never worse than $1/\alpha$.
The extra factor observed in local work comes from the space--time footprint
of repeated scans.

\item \emph{Peak-volume flattening is rigorously feasible.}
RPPR KKT structure supplies a locally checkable, momentum-independent gate
whose working set never exceeds volume $1/\rho$.  This part no longer depends
on a conjectural signed-residual $\ell_1$ invariant.

\item \emph{Naive restart and batching mechanisms are insufficient.}
Short fixed restart blocks lose acceleration; FIFO batch splitting preserves
triangular work; and a path forces accelerated-scale many safe expansions.
Even one endpoint edge forbids carrying a zero fixed-face energy through an
admission without an explicit optimum-shift term.

\item \emph{The open problem is now narrow.}
Exact center transport removes the face-change defect and has a fully charged
new-volume implementation on paths.  Weighted shocks are now exactly an
occupancy charge on the remaining face gain, and the unweighted Schur
telescope alone cannot pay that charge with a factor $\sqrt\alpha$.  One must
still pack the zero-padded cross term or control this occupancy jointly with
the stage count and number of old-prefix sweeps, then
extend the path response beyond one-dimensional append-only geometry.  Until
then, the graph-uniform
$\softO(1/(\sqrt\alpha\epsppr))$ theorem remains open.
The exact zero-start family in
\cref{prop:path-zero-start-swept-separation} additionally shows that a
log-free realized-volume charge is false for the literal full-prefix
implementation, but its certified lower-bound factor is only an accuracy
logarithm and hence does not close this soft-order question.  At the constant
ratio, \cref{prop:path-constant-ratio-floor} fixes
$J,\nu_{\rm fin}=\Theta(q^{-1})$ and only the product-scale lower floor
$\mathfrak V_T=\Omega(\nu_{\rm fin}/q)$; the stage-$6$ obstruction in
\cref{prop:path-frontier-only-gate-fails} shows why the missing chronology
cannot be reduced to the raw frontier coordinate.  The moving-maximum bound
in \cref{thm:path-moving-max-soft-upper} now proves
$T=O(q^{-2}\log(1/q))$ and
$\mathfrak V_T=O(q^{-3}\log(1/q))$ without fixing the maximizer.  It leaves
the matching global lower powers and aggregate logarithm removal open and
therefore does not establish the sampled $\Theta$ laws.  The audited short
full-path candidate proves only the conditional roots in
\cref{lem:path-full-face-linearized-roots} and three finite exact traces; the
face-uniform $O(q^{-1})$ route and a logarithmic fixed-face lower bound are
both open.  The exact stage-$6$--$7$ STOP in
\cref{prop:path-monotone-correction-potential-fails} also shows that the next
upper attempt cannot simply declare the moving correction to be decreasing,
or bound it by $q^{-1}$ times the current normalized face error with unit
coefficient.  The stage-$9$--$10$ STOP in
\cref{prop:path-correction-error-bank-potential-fails} further shows that
adding that same coefficient-one face error as a decreasing bank does not
repair the argument, even though the bank passes the earlier correction
spike.  The exact inequalities in
\cref{prop:path-no-constant-coefficient-bank} strengthen this to every
nonnegative constant coefficient: satisfying the local $U_3$ pair violates
the local $U_4$ pair, or conversely.  They do not analyze a multi-step block
or a global block horizon.  The tight endpoint coefficients do give the
squared, Schur-paid isolated reset in
\cref{prop:path-tight-pair-schur-bank}.  The exact extension in
\cref{prop:path-tight-pair-local-chain-stop} proves why that charge does not
turn the surrounding fixed-face evolution into a pointwise potential: the
production step fails first, the reset passes, the follow-up fails, and their
complete-chain net change is positive.  A later fixed-face recovery does
eventually pay that finite deficit: the telescope in
\cref{prop:path-tight-pair-recovery-block} stops through stage $12$ and first
closes at stage $13$.  Its rectangle-uniform endpoint inequality is one
local block fact and gives no uniform face horizon or repeated-admission
potential.
The face-general causal identity gives a cleaner scalar interface, but its two
finite tests split sharply.  Carried credit keeps the canonical path solvent
across two admissions, whereas the asymmetric T tree retains negative debt
through its next admission and only recovers later.  Thus any surviving
global use cannot rely only on a uniform restarted-block horizon.  More
strongly, \cref{prop:three-admission-all-history-stop} retains all history
from the actual zero initialization and still becomes negative after three
consecutive singleton admissions.  Hence the unqualified zero-balance
$\delta^2$ all-history route is closed for the named recurrence and gate.  A
surviving argument must impose a structural condition, justify extra initial
reserve, or use a different observable or energy account.
The consumed-energy account in
\cref{prop:consumed-energy-structural-solvency} supplies one such causal
reserve without adding initial balance: it stores only energy already lost on
executed fixed-face contractions.  It repairs the six-vertex STOP and the
finite rooted graph atlas at $q=1/5$ with small coefficients.  The leaf-seeded
star in \cref{prop:consumed-energy-star-stop} nevertheless forces coefficient
$\Omega(q^{-3})$, so no constant or polylogarithmic multiplier yields the
desired uniform repair.  On interior traces of the named complete gate,
admission preserves the score exactly, and the baseline energy comparison is
solvent with $O(q^{-5})$.  The retained projection-normal identity improves
this to $(C_{\rm n}+1-q^2)/q^4$ under the causal normal--anchor condition and
to $(1-q^2)/q^4$ when projection is inactive.  The support-aware safe envelope in
\cref{prop:consumed-reserve-quartic-projection} instead gives
$(1-q^2)/q^4$ under arbitrary projection; its score ignores clipped-zero
rows, but its admission and terminal scans do not.  The reachable
$K_{2,3}$-plus-leaf family in
\cref{prop:consumed-energy-quartic-lower} remains projection-inactive through
its reviewed stage and requires coefficient
$49/(3884q^4)+O(q^{-3})$.  Hence the support-aware and inactive-original
coefficient order is sharply $\Theta(q^{-4})$.  In particular, the explicit
low-frequency condition cannot hold with graph-independent constant on all
reachable traces, though it may define a useful narrower promised class.  The
first-admission reserve floor in
\cref{thm:original-score-projected-quartic} closes that last scalar gap.  A
singleton boundary admission forces $E_0=\Omega(q^5)$ and hence already
consumed reserve $R=\Omega(q^6)$; this pays the universal $q^2/25$ ceiling on
every clipped-zero score.  The original all-coordinate score therefore has
the same $(1-q^2)/q^4$ coefficient under arbitrary projection, and its
graph-uniform order is also sharply $\Theta(q^{-4})$.  The threshold-tuned
endpoint path and the reachable 30-vertex clipped row remain useful
mechanism tests, but no longer define a scalar-order gap.  None of these balance statements
alone
controls convergence or total work, and every uncached optimum, energy,
support, or gate evaluation must be charged.
\end{enumerate}

The gate is compatible with the broader hybrid philosophy: use signed AESP,
Chebyshev, SOR, or another accelerator to generate candidates, but let a
monotone KKT envelope decide which coordinates may be exposed.  A
momentum-free local cleanup remains a proof-safe fallback.  The new result
does not close the separate early-AESP locality promotion gate in the
repository's AESP--\LOCSOR{} note; it instead supplies an alternative route
whose own continuation lemma must be resolved.
